\section{Návrh reprodukovateľného experimentu}
\label{sec:methodology}

Experiment je navrhnutý tak, aby bolo možné oddelene skontrolovať
každý krok: výber záznamov, získanie verejnej transakcie, výpočet
príznakov, vytvorenie grafu, rozdelenie dát a trénovanie modelu.
Zdrojový kód sa nachádza v adresári \texttt{analysis/} a poskytuje
príkazy \texttt{prepare}, \texttt{collect}, \texttt{mapping-eval},
\texttt{build-datasets}, \texttt{train} a \texttt{report}.

\subsection{Dátový tok}

Obrázok~\ref{fig:pipeline} ukazuje celý postup. Pôvodné súbory Elliptic
a verejné mapovanie sa najprv obmedzia na vybrané časové kroky. Klient
Blockbook potom získa verejné transakcie a uloží každú úspešnú odpoveď
do databázy SQLite. Z odpovedí sa vypočítajú vysvetliteľné príznaky a
hrany. Posledné dve fázy vytvoria modelové varianty a súhrnné tabuľky
a grafy.

\begin{figure}[htbp]
  \centering
  \resizebox{\textwidth}{!}{%
  \begin{tikzpicture}[
      >=Latex,
      node distance=7mm,
      every node/.style={font=\small,align=center},
      processbox/.style={draw,rounded corners,minimum width=47mm,minimum height=9mm},
      data/.style={processbox,fill=blue!8},
      action/.style={processbox,fill=orange!12}
    ]
    \node[data] (elliptic) {Elliptic + verejné mapovanie};
    \node[action,below=of elliptic] (prepare) {\texttt{prepare}: výber 6 časových krokov};
    \node[action,below=of prepare] (collect) {\texttt{collect}: Blockbook + pamäť SQLite};
    \node[action,below=of collect] (features) {výpočet 18 lokálnych príznakov a hrán};
    \node[action,below left=8mm and 6mm of features] (mapping) {\texttt{mapping-eval}\\skrytých 20~\%};
    \node[action,below right=8mm and 6mm of features] (datasets) {\texttt{build-datasets}\\3 varianty, 2 rozdelenia};
    \node[action,below=18mm of features] (train) {\texttt{train}: LR + RF + GCN};
    \node[data,below=of train] (summary) {\texttt{report}: tabuľky, grafy, manifest};
    \draw[->] (elliptic) -- (prepare);
    \draw[->] (prepare) -- (collect);
    \draw[->] (collect) -- (features);
    \draw[->] (features) -- (mapping);
    \draw[->] (features) -- (datasets);
    \draw[->] (datasets) -- (train);
    \draw[->] (mapping.south) -- (summary.west);
    \draw[->] (train) -- (summary);
  \end{tikzpicture}
  }
  \caption{Dátový tok experimentu. Vlastné spracovanie.}
  \label{fig:pipeline}
\end{figure}

Každý beh vytvorí manifest s cestami a kryptografickými odtlačkami
vstupov, verziou Pythonu a knižníc, počtami záznamov, vybranými krokmi,
rozsahmi blokov a počiatočnou hodnotou generátora náhodných čísel
(angl. \emph{seed}). Surové súbory, mapovanie, vyrovnávacia pamäť API
a odvodený dátový súbor sa neukladajú do systému Git. Dôvodom je veľkosť
údajov aj licencia CC BY-NC-ND 4.0, ktorá obmedzuje ďalšie šírenie
upravených verzií Elliptic \cite{safar2026elliptic,ellipticDataset}.

Tabuľka~\ref{tab:own-pipeline-commands} spresňuje, čo každý príkaz
načíta a čo po sebe zanechá. Takéto oddelenie umožňuje zopakovať iba
zlyhaný krok a súčasne skontrolovať jeho medzivýsledok.

\begin{table}[p]
  \centering
  \small
  \caption{Príkazy vlastného experimentálneho postupu. Vlastné spracovanie.}
  \label{tab:own-pipeline-commands}
  \begin{tabularx}{\textwidth}{lXX}
    \toprule
    Príkaz & Čo vykoná & Kontrolovateľný výstup \\
    \midrule
    \texttt{prepare} & načíta tri súbory Elliptic a verejné mapovanie,
      vyberie šesť krokov a overí identifikátory & vybrané príznaky,
      triedy, hrany, mapovanie a manifest \\
    \texttt{collect} & určí bloky vybraných krokov a cez Blockbook
      načíta všetky kandidátske transakcie & SQLite cache a tabuľka
      transakcií s 18 vypočítanými hodnotami \\
    \texttt{mapping-eval} & skryje 20~\% známych zhôd v každom kroku,
      kalibruje porovnanie iba na zostávajúcich 80~\% & top-1, top-5,
      nejednoznačné a nenájdené prípady \\
    \texttt{build-datasets} & vytvorí tri kombinácie príznakov a grafu,
      časové a adresne oddelené rozdelenie & vrcholy, dve sady hrán a
      masky tréningu, validácie a testu \\
    \texttt{train} & normalizuje podľa tréningu, vyberie hyperparametre
      na validácii a vyhodnotí tri modely & každý beh, metriky, prah a
      agregovaný priemer so smerodajnou odchýlkou \\
    \texttt{report} & číta uložené výsledky bez nového učenia
      & grafy PDF/PNG použité v práci \\
    \bottomrule
  \end{tabularx}
\end{table}

\subsection{Výber časových krokov}

Výber obsahuje dva trénovacie, dva validačné a dva testovacie časové
kroky. Umožňuje zopakovať celý postup na 28~221 vrcholoch bez potreby
spracovať všetkých 203~769 transakcií. Konkrétne počty uvádza
tabuľka~\ref{tab:selected-steps}.

\begin{table}[htbp]
  \centering
  \caption{Vybrané časové kroky a rozdelenie experimentu. Vlastné
  spracovanie z pôvodných súborov Elliptic.}
  \label{tab:selected-steps}
  \begin{tabular}{lrrrrr}
    \toprule
    Časový krok & Časť & Všetky & Licitné & Ilicitné & Neoznačené \\
    \midrule
    20 & tréning & 4~291 & 640 & 260 & 3~391 \\
    29 & tréning & 4~275 & 845 & 329 & 3~101 \\
    30 & validácia & 2~483 & 441 & 83 & 1~959 \\
    32 & validácia & 4~525 & 981 & 342 & 3~202 \\
    35 & test & 5~507 & 1~159 & 182 & 4~166 \\
    42 & test & 7~140 & 1~915 & 239 & 4~986 \\
    \midrule
    Spolu & & 28~221 & 5~981 & 1~435 & 20~805 \\
    \bottomrule
  \end{tabular}
\end{table}

Na obrázku~\ref{fig:selected-split-timeline} je základnou jednotkou
rozdelenia celý časový krok. Vrcholy v rámci kroku sa náhodne nedelia.
Kroky 20 a 29 sú
jediné dáta, z ktorých model upravuje svoje parametre. Kroky 30 a 32
slúžia na výber nastavení a rozhodovacieho prahu. Kroky 35 a 42 sa
použijú až pri konečnom meraní. Medzery medzi číslami pripomínajú, že
výber nepokrýva súvislé obdobie blockchainu.

\begin{figure}[htbp]
  \centering
  \resizebox{\textwidth}{!}{%
  \begin{tikzpicture}[
      >=Latex,
      every node/.style={font=\small,align=center},
      cell/.style={draw,rounded corners,minimum width=23mm,minimum height=14mm}
    ]
    \node[cell,fill=blue!12] (s20) at (0,0) {krok 20\\4~291};
    \node[cell,fill=blue!12] (s29) at (2.7,0) {krok 29\\4~275};
    \node[cell,fill=yellow!20] (s30) at (5.4,0) {krok 30\\2~483};
    \node[cell,fill=yellow!20] (s32) at (8.1,0) {krok 32\\4~525};
    \node[cell,fill=orange!18] (s35) at (10.8,0) {krok 35\\5~507};
    \node[cell,fill=orange!18] (s42) at (13.5,0) {krok 42\\7~140};
    \draw[->,thick] (-1.4,-1.1) -- (14.9,-1.1) node[right]{čas};
    \draw[decorate,decoration={brace,amplitude=5pt,mirror}]
      (-1.1,0.95) -- (3.8,0.95) node[midway,yshift=9mm]{tréning};
    \draw[decorate,decoration={brace,amplitude=5pt,mirror}]
      (4.3,0.95) -- (9.2,0.95) node[midway,yshift=9mm]{validácia};
    \draw[decorate,decoration={brace,amplitude=5pt,mirror}]
      (9.7,0.95) -- (14.6,0.95) node[midway,yshift=9mm]{test};
  \end{tikzpicture}%
  }
  \caption{Časové rozdelenie 28~221 vybraných vrcholov. Číslo pod
  krokom udáva počet všetkých transakcií v danom kroku. Vlastné
  spracovanie z Elliptic.}
  \label{fig:selected-split-timeline}
\end{figure}

Verejné mapovanie obsahuje 28~152 z vybraných 28~221 identifikátorov,
teda 99,76~\% tejto vzorky. Zvyšných 69 záznamov nemá potvrdené
mapovanie. Program pre ne vypíše kandidátov s mierou nejednoznačnosti a
priradí im stav \texttt{candidate\_only}.

\subsection{Zber údajov cez Blockbook}

Blockbook je indexer blockchainu s rozhraním REST. Pre jednu transakciu
vracia okrem \texttt{txid} aj výšku a čas bloku, hodnotu, poplatok,
veľkosť, vstupy a výstupy. Hodnoty sú vyjadrené v najmenšej jednotke
daného blockchainu, pri Bitcoine v satoshi \cite{blockbook}.

Klient používa vlastnú hlavičku \texttt{User-Agent}, časový limit a
exponenciálne predlžovanie prestávky medzi opakovanými pokusmi. Najviac
štyri pracovníky odosielajú požiadavky súbežne, aby sa verejný server
nepreťažoval. Každá odpoveď s kódom 200 sa uloží do SQLite podľa úplnej
adresy požiadavky. Pri opakovanom behu sa teda číta z vyrovnávacej
pamäte a zber môže
pokračovať po prerušení.

Pre vyhodnotenie párovania program prejde všetky transakcie v blokoch
priradených ku konkrétnemu časovému
kroku. Tým vytvorí množinu kandidátov, v ktorej musí algoritmus nájsť
správny verejný \texttt{txid} bez použitia skrytej odpovede.

\subsection{Transparentné lokálne príznaky}

Z jednej odpovede Blockbook sa počíta 18 príznakov so známym významom.
Príznak je číselná hodnota, ktorou sa transakcia opíše modelu.
Samostatné číslo neurčuje triedu; model hľadá kombinácie viacerých
čísel. Tabuľka~\ref{tab:transparent-features} uvádza všetkých 18
hodnôt, ich jednotku a význam.

\begin{table}[p]
  \centering
  \scriptsize
  \caption{Transparentne vypočítané lokálne príznaky.}
  \label{tab:transparent-features}
  \begin{tabularx}{\textwidth}{r l l X}
    \toprule
    Č. & Príznak & Jednotka & Čo model dostane \\
    \midrule
    2 & objem & satoshi & súčet všetkých výstupov transakcie \\
    3 & poplatok & satoshi & vstupná hodnota mínus výstupná hodnota \\
    5 & počet vstupov & počet & koľko starších výstupov sa míňa \\
    6 & počet výstupov & počet & koľko nových výstupov vzniká \\
    7 & jedinečné vstupné adresy & počet & počet rôznych neprázdnych adries vo vstupoch \\
    9 & minimum vstupov & satoshi & najmenší vstup \\
    10 & maximum vstupov & satoshi & najväčší vstup \\
    11 & odchýlka vstupov & satoshi & rozdielnosť veľkostí vstupov \\
    12 & priemer vstupov & satoshi & priemerná hodnota jedného vstupu \\
    13 & Pearson vstupov & $[-1,1]$ & lineárny trend hodnôt podľa indexu \\
    14 & Spearman vstupov & $[-1,1]$ & rast alebo pokles poradia hodnôt \\
    15 & jedinečné výstupné adresy & počet & počet rôznych neprázdnych adries vo výstupoch \\
    17 & minimum výstupov & satoshi & najmenší výstup \\
    18 & maximum výstupov & satoshi & najväčší výstup \\
    19 & odchýlka výstupov & satoshi & rozdielnosť veľkostí výstupov \\
    20 & priemer výstupov & satoshi & priemerná hodnota jedného výstupu \\
    21 & Pearson výstupov & $[-1,1]$ & lineárny trend hodnôt podľa indexu \\
    22 & Spearman výstupov & $[-1,1]$ & rast alebo pokles poradia hodnôt \\
    \bottomrule
  \end{tabularx}
\end{table}

Pri príklade z obrázka~\ref{fig:utxo-example} by počet vstupov bol 2,
počet výstupov 2, minimum vstupov 0,4 BTC, maximum vstupov 0,6 BTC,
priemer vstupov 0,5 BTC a poplatok 0,01 BTC. Počet jedinečných adries
sa môže od počtu vstupov líšiť, pretože jedna adresa sa môže objaviť
viackrát. Program najprv vytvorí množinu adries a až potom ju spočíta.

Pri prázdnej postupnosti, napríklad pri vstupoch coinbase transakcie,
sa agregačná hodnota nastaví na nulu. Smerodajná odchýlka jednej
hodnoty a korelácia konštantnej postupnosti sú tiež nula. Toto pravidlo
je súčasťou implementácie a testov, pretože matematická korelácia v
takých prípadoch nie je určená.

\subsection{Kontrolovaný experiment párovania}

Verejné mapovanie už obsahuje známe odpovede. Presnosť metódy sa preto
meria po skrytí 20~\% známych dvojíc v každom zo šiestich časových
krokov pomocou \texttt{seed=42}. Zvyšných 80~\% slúži na
odhad lineárneho prevodu medzi normalizovaným príznakom Elliptic a
hodnotou vypočítanou z blockchainu:
\begin{equation}
  \hat{x}=a z+b.
\end{equation}
Koeficienty $a$ a $b$ sa určia metódou najmenších štvorcov. Skryté
\texttt{txid} sa nepoužijú pri odhade ani pri tvorbe kandidátov.

Vyhľadávanie sa hodnotí v troch stupňoch:
\begin{enumerate}
  \item iba počet vstupov a výstupov;
  \item počty doplnené o objem, poplatok a počty jedinečných adries;
  \item všetkých 18 obnovených lokálnych príznakov.
\end{enumerate}
Najprv sa kandidáti zúžia na rovnaký počet vstupov a výstupov. Potom
sa zoradia podľa priemernej normalizovanej absolútnej odchýlky ostatných
príznakov. Zaznamenáva sa úspešnosť top-1 a top-5, počet kandidátov,
chýbajúce výsledky a prípady, keď dvaja kandidáti dosiahnu rovnaké
skóre. Takéto prípady sa označia ako nejednoznačné.

\subsection{Varianty dát a grafu}

Experiment porovnáva tri varianty:
\begin{enumerate}[label=\Alph*.]
  \item pôvodné anonymizované príznaky a pôvodný graf Elliptic;
  \item transparentné lokálne príznaky a pôvodný graf Elliptic;
  \item transparentné lokálne príznaky a graf rekonštruovaný z odkazov
  vstupov na staršie \texttt{txid}.
\end{enumerate}

Variant B oddeľuje vplyv príznakov, kým rozdiel medzi B a C oddeľuje
vplyv grafu. Logistická regresia a náhodný les graf nepoužívajú, preto
pri nich varianty B a C slúžia ako kontrola a musia pri rovnakej
počiatočnej hodnote
poskytnúť rovnaký výsledok. Pri GCN sa graf priamo podieľa na výpočte
reprezentácie vrcholu.

Podobný dôraz na vopred určené rozdelenia, spoločný formát vstupu a
jednotný hodnotiaci postup používa aj Open Graph Benchmark. Jeho
autori zdôrazňujú, že rozdelenie má napodobniť konkrétne budúce
použitie modelu. Náhodné percentá riadkov túto požiadavku nespĺňajú
\cite{hu2020ogb}. Pre túto prácu z toho vyplýva, že názov algoritmu bez
opisu času, grafových hrán a spoločných adries nestačí na porovnanie.

\subsection{Časové rozdelenie a rozdelenie odolné voči úniku informácií}

Bežné časové rozdelenie priradí celé kroky 20 a 29 do tréningu, kroky
30 a 32 do validácie a kroky 35 a 42 do testu. Tým sa do rôznych častí
náhodne nerozdelia vrcholy jedného časového komponentu Elliptic. Adresa
\texttt{A} sa napriek tomu môže objaviť v trénovacej
transakcii a znova v testovacej transakcii. Model potom môže využiť
opakovanú entitu.

Druhé rozdelenie sa tvorí postupne zľava doprava. Najprv sa prijmú
trénovacie transakcie a zapamätajú sa všetky ich vstupné a výstupné
adresy. Validačná transakcia sa prijme iba vtedy, keď nepoužíva žiadnu
zapamätanú adresu. Jej adresy sa potom tiež zapamätajú. Rovnaké pravidlo
sa nakoniec použije na test. Tento postup zachováva opakovania adresy
v rámci jednej časti a vylúči jej výskyt v neskoršej časti.

\begin{figure}[htbp]
  \centering
  \begin{tikzpicture}[
      >=Latex,
      node distance=7mm,
      every node/.style={font=\small,align=center},
      box/.style={draw,rounded corners,minimum width=40mm,minimum height=11mm}
    ]
    \node[box,fill=blue!12] (train) {tréning\\ulož adresy $\{A,B\}$};
    \node[box,right=of train,fill=yellow!20] (val) {validácia\\odstráň riadok s $A$};
    \node[box,right=of val,fill=orange!18] (test) {test\\odstráň riadok s $A,B$};
    \draw[->] (train) -- node[above,font=\scriptsize]{zakázané v ďalšej časti} (val);
    \draw[->] (val) -- node[above,font=\scriptsize]{rozšírený zoznam} (test);
    \node[box,below=12mm of val,fill=gray!10,minimum width=48mm]
      (edges) {nakoniec odstráň hrany\\medzi rozdielnymi časťami};
    \draw[->] (train.south) -- (edges.west);
    \draw[->] (val.south) -- (edges.north);
    \draw[->] (test.south) -- (edges.east);
  \end{tikzpicture}
  \caption{Tvorba rozdelenia bez pozorovateľného prekryvu adries.
  Písmená $A$ a $B$ sú ilustračné adresy. Vlastné spracovanie.}
  \label{fig:clean-split-process}
\end{figure}

Po filtrovaní vrcholov sa odstránia aj všetky hrany, ktorých konce
patria do rozdielnych častí. GCN preto nemôže preniesť informáciu z
trénovacieho vrcholu do validačného alebo testovacieho vrcholu.
Transakcie bez overeného verejného mapovania sa do tohto rozdelenia
nezaradia, pretože pri nich nemožno spoľahlivo skontrolovať adresy.
Počty vyradených vrcholov sa uvádzajú spolu s výsledkami; metrika po
filtrovaní totiž vzniká na menšej a odlišnej vzorke.

Pri variante A zostáva časový únik ukrytý v pôvodných adresných
príznakoch 23--94. Rozdelenie odstraňuje pozorovateľný prekryv adries a
hrán medzi časťami. Informácia z budúcnosti už zapísaná v jednom
vektore zostáva zachovaná. Transparentné varianty B a C tieto globálne
adresné príznaky nepoužívajú.

\subsection{Modely a hyperparametre}

Experiment používa tri referenčné algoritmy ako kontrolné body:
lineárny model, nelineárny tabuľkový model a model využívajúci graf.
Práca nepripisuje týmto algoritmom vlastný metodický prínos. Ak bude v
ďalšej verzii práce vytvorený vlastný
klasifikátor, musí sa porovnať na rovnakých vstupoch, rozdeleniach a
metrikách. Číslo prevzaté z experimentu s inými podmienkami neposkytuje
porovnateľný základ.

\begin{table}[htbp]
  \centering
  \small
  \caption{Presné algoritmy použité v implementácii. Vlastné spracovanie.}
  \label{tab:implemented-models}
  \begin{tabularx}{\textwidth}{lXX}
    \toprule
    Model & Implementácia & Čo sa počas učenia mení \\
    \midrule
    logistická regresia &
      \texttt{LogisticRegression} z knižnice scikit-learn,
      vyvážené váhy tried, najviac 2~000 iterácií &
      jedna váha pre každý príznak a posun $b$; validácia vyberá
      $C\in\{0{,}1;1;10\}$ a rozhodovací prah \\
    náhodný les &
      \texttt{RandomForestClassifier} z knižnice scikit-learn,
      500 stromov &
      otázky v uzloch stromov; validácia vyberá hĺbku 10, 20 alebo bez
      limitu a minimálne 1 alebo 5 vzoriek v liste \\
    dvojvrstvová GCN &
      vlastná trieda v PyTorch, dve lineárne vrstvy 128$\rightarrow$2,
      riedka normalizovaná matica susednosti &
      matice váh oboch vrstiev; Adam používa krok 0,01, váhový útlm
      $5\cdot10^{-4}$, dropout 0,5 a skoré zastavenie po 8 epochách \\
    \bottomrule
  \end{tabularx}
\end{table}

Váhy neurónového a lineárneho modelu sa všeobecne menia proti smeru,
v ktorom rastie stratová funkcia:
\begin{equation}
  \boldsymbol{\theta}_{t+1}
  =\boldsymbol{\theta}_{t}-\eta\nabla_{\boldsymbol{\theta}}
  \mathcal{L}(\boldsymbol{\theta}_{t}),
  \label{eq:gradient-update}
\end{equation}
kde $\boldsymbol{\theta}$ označuje učené parametre, $\eta$ rýchlosť
učenia a $\nabla\mathcal{L}$ smer najväčšieho rastu chyby. V praxi GCN
používa adaptívny algoritmus Adam, ktorý pre každý parameter upravuje
veľkosť kroku podľa priebehu predchádzajúcich gradientov. Náhodný les
gradient nepoužíva; skladá rozhodovacie pravidlá priamo z trénovacích
príkladov.

Pred učením sa každý príznak prevedie na porovnateľnú mierku. Ak sú
napríklad poplatky v tisícoch satoshi a korelácia v intervale
$[-1,1]$, veľké jednotky poplatku by mohli neprimerane ovplyvniť
optimalizáciu. Štandardizovaná hodnota je
\begin{equation}
  z_{ij}=\frac{x_{ij}-\mu_j^{\mathrm{train}}}
  {\sigma_j^{\mathrm{train}}},
  \label{eq:training-standardization}
\end{equation}
kde $x_{ij}$ je príznak $j$ transakcie $i$. Priemer $\mu_j$ a
smerodajná odchýlka $\sigma_j$ sa počítajú výhradne z tréningu.
Validácia a test používajú tie isté trénovacie hodnoty. Model tak pri
príprave vstupu nevidí rozdelenie budúcich dát.

\subsubsection{Logistická regresia}

Logistická regresia priradí každému príznaku váhu. Kladná váha zvyšuje
skóre ilicitnej triedy, záporná ho
znižuje a váha blízka nule znamená malý lineárny vplyv. Pre vektor
príznakov $\mathbf{x}$ najprv vypočíta
\begin{equation}
  s=\mathbf{w}^{\mathsf T}\mathbf{x}+b,
\end{equation}
kde $\mathbf{w}$ sú učené váhy a $b$ je posun. Sigmoidálna funkcia
prevedie neobmedzené skóre na číslo od 0 do 1:
\begin{equation}
  p(y=1\mid\mathbf{x})=\frac{1}{1+\mathrm{e}^{-s}}.
  \label{eq:logistic-sigmoid}
\end{equation}
Ak je napríklad výsledok 0,82, model priradil transakcii skóre 0,82
pre ilicitnú triedu. Skóre je výstup modelu kalibrovaného voči
označeniam Elliptic a nemá význam právnej pravdepodobnosti.

Počas učenia sa váhy menia tak, aby znižovali binárnu krížovú entropiu
\begin{equation}
  \mathcal{L}=-\frac{1}{n}\sum_{i=1}^{n}
  \left[y_i\log p_i+(1-y_i)\log(1-p_i)\right].
\end{equation}
Sebavedomá chybná predikcia vytvorí veľkú stratu. Regularizácia bráni
neprimerane veľkým váham. Implementácia porovnáva
$C\in\{0{,}1;1;10\}$ a vyberá hodnotu podľa validačného $F_1$.
Váhy tried sú vyvážené, aby početnejšia licitná trieda nepotlačila
menšiu ilicitnú triedu.

\subsubsection{Náhodný les}

Rozhodovací strom sa učí sériu otázok typu \uv{je poplatok menší než
hranica $t$?} alebo \uv{je počet výstupov väčší než 2?}. Každá otázka
rozdelí trénovacie príklady na dve vetvy. Algoritmus vyberá príznak a
hranicu, ktoré znižujú nečistotu tried. Giniho nečistota uzla je
\begin{equation}
  G=1-\sum_{k\in\{0,1\}}p_k^2,
  \label{eq:gini}
\end{equation}
kde $p_k$ je podiel triedy $k$ v uzle. Ak uzol obsahuje iba jednu
triedu, $G=0$. Pri rovnakom zastúpení dvoch tried je $G=0{,}5$.

Jeden strom sa môže príliš prispôsobiť náhodným detailom tréningu.
Náhodný les preto vytvorí 500 stromov. Každý strom dostane náhodnú
vzorku trénovacích riadkov a pri delení skúma náhodnú podmnožinu
príznakov. Výsledné skóre je priemer hlasov stromov. Model tak dokáže
zachytiť nelineárne pravidlá, napríklad kombináciu nízkeho počtu
výstupov, vysokej hodnoty a nezvyčajného poplatku
\cite{breiman2001random}.

Ak $p_m(y=1\mid\mathbf{x})$ je podiel ilicitných trénovacích príkladov
v liste stromu $m$, skóre celého lesa s $M=500$ stromami je
\begin{equation}
  p_{\mathrm{RF}}(y=1\mid\mathbf{x})
  =\frac{1}{M}\sum_{m=1}^{M}p_m(y=1\mid\mathbf{x}).
  \label{eq:rf-average}
\end{equation}
Les teda nehľadá jednu jedinú charakteristiku \uv{ilicitnosti}. Každý
strom vytvorí inú postupnosť otázok a konečné skóre ich spriemeruje.

Porovnávajú sa maximálne hĺbky 10, 20 a bez obmedzenia a minimálne
počty vzoriek v liste 1 alebo 5. Najlepšia kombinácia sa zvolí na
validácii pri počiatočnej hodnote 42. Potom sa celý les naučí päťkrát
s hodnotami 42 až 46. Rozdiel medzi behmi ukazuje, ako výsledok závisí
od náhodného výberu vzoriek a príznakov.

\subsubsection{Grafová konvolučná sieť}

Logistická regresia a náhodný les posudzujú každý riadok samostatne.
Grafová konvolučná sieť (GCN) doplní k príznakom transakcie aj
informáciu od jej susedov. Pred výpočtom sa ku každému vrcholu pridá
hrana do seba a matica susednosti sa normalizuje:
\begin{equation}
  \hat{A}=A+I,\qquad
  \tilde{A}=\hat{D}^{-1/2}\hat{A}\hat{D}^{-1/2}.
\end{equation}
Matica $A$ opisuje hrany, $I$ pridáva vlastnú informáciu vrcholu a
$\hat{D}$ obsahuje počty spojení. Normalizácia bráni tomu, aby vrchol s
mnohými susedmi získal automaticky neúmerne veľké hodnoty.

Jedna grafová vrstva sa dá zapísať ako
\begin{equation}
  H^{(l+1)}=\sigma\!\left(\tilde{A}H^{(l)}W^{(l)}\right),
  \label{eq:gcn-layer}
\end{equation}
kde $H^{(0)}=X$ je vstupná matica príznakov, $W^{(l)}$ sú učené váhy a
$\sigma$ je aktivačná funkcia. Prvá vrstva v implementácii vytvorí pre
každý vrchol 128 skrytých hodnôt. Použitá funkcia ReLU ponechá kladné
číslo a záporné nahradí nulou:
\begin{equation}
  \operatorname{ReLU}(z)=\max(0,z).
  \label{eq:relu}
\end{equation}
Nelinearita umožňuje dvom vrstvám modelovať vzťahy, ktoré jedna
lineárna transformácia nezachytí. Po ReLU sa počas tréningu náhodne vynuluje polovica
skrytých hodnôt. Tento postup sa nazýva náhodné vypínanie, angl.
\emph{dropout}. Model sa tým učí používať viac skrytých hodnôt.

Druhá grafová vrstva vytvorí dve skóre, jedno pre licitnú a druhé pre
ilicitnú triedu. Funkcia softmax ich prevedie na súčet 1:
\begin{equation}
  p_k=\frac{\mathrm{e}^{s_k}}
  {\mathrm{e}^{s_0}+\mathrm{e}^{s_1}},\qquad k\in\{0,1\}.
  \label{eq:softmax}
\end{equation}
Váhy sa
upravujú algoritmom Adam s rýchlosťou učenia 0,01 a váhovým útlmom
$5\cdot10^{-4}$. Používa sa vážená krížová entropia, aby ilicitná
trieda mala napriek menšiemu počtu príkladov primeraný vplyv. Tréning
má najviac 200 epoch; epocha je jeden úplný krok učenia nad
trénovacími vrcholmi. Ak sa validačný $F_1$ nezlepší osem epoch, uloží
sa najlepší skorší stav a učenie sa zastaví \cite{kipf2017gcn}.

Dve grafové vrstvy znamenajú, že reprezentácia vrcholu môže obsahovať
informáciu od susedov vzdialených najviac dva preskoky. Neoznačené
vrcholy zostávajú v grafe a môžu sprostredkovať túto informáciu. Keďže
nemajú cieľovú triedu, nevstupujú do stratovej funkcie ani do metrík.

\subsubsection{Spoločný postup učenia}

\begin{figure}[htbp]
  \centering
  \resizebox{\textwidth}{!}{%
  \begin{tikzpicture}[
      >=Latex,
      node distance=6mm,
      every node/.style={font=\small,align=center},
      flowbox/.style={draw,rounded corners,minimum width=31mm,minimum height=13mm}
    ]
    \node[flowbox,fill=blue!10] (train) {tréning\\uprav váhy};
    \node[flowbox,right=of train,fill=yellow!18] (validation) {validácia\\vyber nastavenia a prah};
    \node[flowbox,right=of validation,fill=green!10] (freeze) {uzamkni model\\už nič neuč};
    \node[flowbox,right=of freeze,fill=orange!16] (test) {test\\zmeraj výsledok raz};
    \draw[->] (train) -- (validation);
    \draw[->] (validation) edge[bend left=22] node[above,font=\scriptsize]{ďalšie nastavenie} (train);
    \draw[->] (validation) -- (freeze);
    \draw[->] (freeze) -- (test);
  \end{tikzpicture}%
  }
  \caption{Oddelenie učenia, výberu nastavení a konečného testu.
  Testovacie dáta sa nevracajú späť do učenia. Vlastné spracovanie.}
  \label{fig:model-training-flow}
\end{figure}

Každý model vráti skóre ilicitnej triedy od 0 do 1. Na validácii sa
skúšajú prahy od 0,05 do 0,95 s krokom 0,01 a vyberie sa prah s
najvyšším ilicitným $F_1$. Ak je napríklad prah 0,63, skóre 0,70 sa
zmení na ilicitnú predikciu a skóre 0,40 na licitnú. Vybraný prah sa
bez ďalšej úpravy použije na test.

\subsection{Metriky a interpretačné pravidlá}

Predikcie sa najprv rozdelia do štyroch skupín. \emph{TP} je správne
zachytená ilicitná transakcia. \emph{FP} je licitná transakcia, ktorú
model omylom označil ako ilicitnú. \emph{FN} je nezachytená ilicitná
transakcia a \emph{TN} je správne rozpoznaná licitná transakcia. Tieto
štyri počty tvoria maticu zámen.

Precíznosť odpovedá na otázku: \uv{Ak model upozornil na transakciu,
ako často mal podľa dát pravdu?}
\begin{equation}
  \mathrm{precision}=\frac{TP}{TP+FP}.
\end{equation}
Úplnosť odpovedá na inú otázku: \uv{Koľko zo všetkých ilicitných
transakcií model našiel?}
\begin{equation}
  \mathrm{recall}=\frac{TP}{TP+FN}.
\end{equation}
Označenie väčšieho počtu transakcií môže zvýšiť úplnosť a súčasne
vytvoriť viac falošných upozornení. Hlavná metrika $F_1$ je
harmonický priemer oboch hodnôt:
\begin{equation}
  F_1 = 2\frac{\mathrm{precision}\cdot\mathrm{recall}}
  {\mathrm{precision}+\mathrm{recall}}.
\end{equation}

Ak je napríklad $TP=80$, $FP=20$ a $FN=40$, precíznosť je
$80/(80+20)=0{,}80$, úplnosť $80/(80+40)=0{,}67$ a $F_1$ približne
0,73. Vysoká obyčajná presnosť by bola pri nevyvážených triedach
zavádzajúca: model by mohol označiť všetko ako licitné a napriek tomu
správne zaradiť väčšinu riadkov.

Vyvážená úspešnosť je priemer správne zachyteného podielu licitnej a
ilicitnej triedy. PR-AUC je plocha pod krivkou precíznosti a úplnosti
pri všetkých možných prahoch. Hodnotí poradie skóre bez výberu jediného
prahu. Hodnota smerujúca k 1 je
lepšia; jej základná úroveň závisí od podielu ilicitnej triedy.

Výsledok modelu sa interpretuje ako predikcia triedy v dátovom súbore
a neposkytuje dôkaz trestného činu. Kandidát párovania vyjadruje
technickú podobnosť s verejnou transakciou bez identifikácie osoby.
